\subsection{A coordinate-level correspondence with the actual Connes quotient}
\label{ccb:bridge}
We now construct the comparison with the actual object of
\cite[Section 3.6, Proposition 3.6]{ccmProlate2024}.
Write $\mathcal H=\mathcal H_{\leq1}$ for its ring of entire functions
of Hadamard order at most one, retaining its stated Hadamard topology.
Write
\[
 N_{\rm C}=\overline{\C[s]\Xi}^{\mathcal H},\qquad
 Q_{\rm C}=\mathcal H/N_{\rm C},\qquad Q_t=\A/N_t.
\]
Here $s$ is the same coordinate in both quotients. The closure defining
$N_{\rm C}$ is precisely the closure of the generators in that proposition;
the following calculation proves the equality of their algebraic spans,
including the transform conventions. The comparison constructed below uses
the actual topology on $\mathcal H$, through a subspace of a product; it does
not substitute the topology of $\A$ for it.

\paragraph{The original summation, Mellin transform, and their coordinates.}
Let $\mathcal G_0$ be the vector space of the finite sums
\begin{equation}
 f(x)=e^{-\pi x^2}\sum_{j=0}^{n}c_j\pi^j x^{2j},\qquad
 c_0=0,\qquad
 \sum_{j=0}^{n}c_j\frac{(2j)!}{4^j j!}=0.
 \label{ccb:gaussian-space}
\end{equation}
The two constraints say exactly $f(0)=\widehat f(0)=0$ for
$\widehat f(y)=\int_\R f(x)e^{-2\pi ixy}dx$. Indeed the Gaussian
moment is $\int_\R\pi^j x^{2j}e^{-\pi x^2}dx=(2j)!/(4^j j!)$.
The original maps of \cite[equations (3.3), (3.4), and Section 3.6]{ccmProlate2024}
are
\begin{align}
 (\mathcal E_{\rm C}f)(\lambda)
   &=\lambda^{1/2}\sum_{n\geq1}f(n\lambda),\qquad \lambda>0,\notag\\
 (\mathcal F_\mu g)(s)
   &=\int_0^\infty g(\lambda)\lambda^{-is}\frac{d\lambda}{\lambda}.
 \label{ccb:original-transforms}
\end{align}
Both signs and the measure are retained. With $\lambda=e^v$, define
\begin{equation}
 k_f(v)=e^{v/2}\sum_{n\geq1}f(ne^v),\qquad
 \iota f(v)=\frac14 k_f(-v).
 \label{ccb:iota}
\end{equation}
Then there is an exact typed map
\begin{equation}
 \iota:\mathcal G_0\longrightarrow\E,\qquad
 T\iota f(s)=\int_\R k_f(v)e^{-isv}dv
             =\mathcal F_\mu\mathcal E_{\rm C}f(s).
 \label{ccb:fourier-bridge}
\end{equation}
The reflection and the factor $1/4$ in \eqref{ccb:iota} are necessary for
the defined $T$, and have not been suppressed.

Here is the convergence proof. For $v\geq0$, every derivative of $k_f$
is bounded by $C_j e^{C_jv}e^{-\pi e^{2v}}$, by differentiating the
finite polynomial times Gaussian and summing the resulting convergent
series in $n$. The Fourier transform of a polynomial times this Gaussian
is again a polynomial times the same Gaussian. Poisson summation and
the two zero constraints give
\[
 \sum_{n\geq1}f(n\lambda)
  =\lambda^{-1}\sum_{n\geq1}\widehat f(n/\lambda),\qquad
 k_f(v)=k_{\widehat f}(-v).
\]
The positive-half-line estimate for $\widehat f$ therefore proves the
estimate on the negative half-line as well, for every derivative.
Every defining seminorm of $\E$ is finite. These bounds justify the
change of variable, reflection, and all entire continuations in
\eqref{ccb:fourier-bridge}.

Put $z=1/2-is$, whose inverse is $s=i(z-1/2)$, and put
\[
 P_f(z)=\sum_{j=0}^n c_j(z/2)_j,\qquad
 (z/2)_j=\prod_{k=0}^{j-1}(z/2+k),\qquad (z/2)_0=1.
\]
For $\Re z>1$, the absolutely convergent Mellin integral gives
\begin{align}
 \mathcal F_\mu\mathcal E_{\rm C}f(s)
 &=\zeta(z)\int_0^\infty f(x)x^{z-1}dx\notag\\
 &=\frac12\pi^{-z/2}\Gamma(z/2)\zeta(z)P_f(z).
 \label{ccb:mellin-polynomial}
\end{align}
The two constraints in \eqref{ccb:gaussian-space} are respectively
$P_f(0)=0$ and $P_f(1)=0$, since $(1/2)_j=(2j)!/(4^j j!)$.
There is consequently a unique polynomial
\begin{equation}
 q_f(s)=\left.\frac{P_f(z)}{z(z-1)}\right|_{z=1/2-is},\qquad
 T\iota f(s)=q_f(s)\Xi(s).
 \label{ccb:qf}
\end{equation}
For the last equality, the right side of \eqref{ccb:mellin-polynomial}
is $q_f(s)\xi(1/2-is)=q_f(s)\Xi(-s)$, and the proved evenness
$\Xi(-s)=\Xi(s)$ supplies the final sign conversion. Entire continuation
then proves the identity on all $\C_s$.

This map is onto every polynomial multiple, with an explicit inverse.
For $q\in\C[s]$, expand the unchanged polynomial
\begin{equation}
 z(z-1)q\bigl(i(z-1/2)\bigr)=\sum_{j=0}^{\deg q+2}c_j(z/2)_j.
 \label{ccb:triangular-inverse}
\end{equation}
There is a unique such expansion because the $j$th basis polynomial
has degree $j$ and leading coefficient $2^{-j}$. Evaluation at $0$
and $1$ proves the two required constraints, giving $f\in\mathcal G_0$.
This proves both injectivity and surjectivity of
$f\mapsto q_f$ and of $f\mapsto q_f\Xi$, since $\Xi$ is not zero
identically. Here is the full basis argument connecting to the source's
Hermite presentation. Its functions are
\[
h_n(x)=\frac{2^{1/4}}{\sqrt{2^n n!}}
 H_n(\sqrt{2\pi}x)e^{-\pi x^2},\qquad
H_n(y)=n!\sum_{j=0}^{\lfloor n/2\rfloor}
 \frac{(-1)^j(2y)^{n-2j}}{j!(n-2j)!}.
\]
Their even members are a triangular basis of all even polynomial
Gaussians, since each has nonzero leading coefficient in its successive
even degree. Put $L=\sqrt\pi x-(2\sqrt\pi)^{-1}\partial_x$.
Differentiating the displayed polynomial proves
$h_n=L^nh_0/\sqrt{n!}$. The Gaussian transform gives
$\widehat h_0=h_0$, while integration by parts gives
$\mathcal F_{\R}L=-iL\mathcal F_{\R}$. Thus
$\widehat h_{2n}=(-1)^nh_{2n}$, with both signs proved in the
original Fourier convention. The finite polynomial formula also gives
$h_{2n}(0)=(-1)^n2^{1/4-n}\sqrt{(2n)!}/n!\ne0$.
Split an even polynomial Gaussian into the spans of the $h_{4\ell}$
and the $h_{4\ell+2}$. The two constraints $f(0)=\widehat f(0)=0$
are exactly vanishing at zero of each of those two components.
In the first span subtracting its constant-basis coordinate gives
$h_{4\ell}-h_{4\ell}(0)h_0/h_0(0)$; in the second it gives
$-h_{4\ell+2}+h_{4\ell+2}(0)h_2/h_2(0)$, with the displayed
minus sign retained. These are precisely the source's
$\psi_\ell^+,\psi_\ell^-$, for $\ell\ge1$. They span the two
constraint kernels by solving for their respective $h_0,h_2$
coefficients. Thus their union spans $\mathcal G_0$, and
\eqref{ccb:qf} proves the required equality of spans.
For the particular original kernel
$h(x)=\pi^2x^4e^{-\pi x^2}-(3\pi/2)x^2e^{-\pi x^2}$,
$P_h(z)=z(z-1)/4$, so $q_h=1/4$, $\iota h=k/4$, and
\begin{equation}
 \mathcal F_\mu\mathcal E_{\rm C}h=\Xi/4,
 \qquad T\iota(4h)=\Xi.
 \label{ccb:original-kernel-factor}
\end{equation}
The factor $4$ in the actual arithmetic function is thereby accounted for.

The original scaling differential operator is
$S=-i(x\partial_x+1/2):\mathcal G_0\to\mathcal G_0$.
It preserves the two constraints: its value at zero is $-if(0)/2$,
and integration by parts gives $\int Sf=(i/2)\int f$.
Termwise differentiation proves
$\mathcal E_{\rm C}Sf=-i\lambda\partial_\lambda\mathcal E_{\rm C}f$.
Integration by parts, with the proved two-ended bounds, then gives
\begin{equation}
 T\iota Sf=M_s T\iota f,\qquad
 \iota Sf=i\partial_v\iota f.
 \label{ccb:scaling-intertwiner}
\end{equation}
The second identity also follows from differentiating $k_f(-v)/4$;
it records the sign produced by the reflection. Thus the comparison
intertwines the actual scaling operator and multiplication by the same $s$.

\paragraph{The correspondence between the two completed quotients.}
Define a space of common representatives, with its topology specified by
the displayed embedding into the actual product:
\begin{equation}
 C=\A\cap\mathcal H,\qquad
 C\longrightarrow\A\times\mathcal H,\quad f\longmapsto(f,f).
 \label{ccb:common-space}
\end{equation}
Give $C$ the subspace topology of this embedding and put
\begin{equation}
 K=N_0\cap N_{\rm C},\qquad \mathscr R=C/K.
 \label{ccb:relation-space}
\end{equation}
Both intersections mean equality of entire functions in the unchanged
$s$ coordinate. In particular $K\subset C$. Since the two quotient
maps are continuous, the following is a continuous injective linear map,
with its complete coordinate formula:
\begin{equation}
 j:\mathscr R\longrightarrow Q_0\times Q_{\rm C},\qquad
 j([f]_K)=([f]_{N_0},[f]_{N_{\rm C}}).
 \label{ccb:quotient-correspondence}
\end{equation}
Indeed the map on $C$ is continuous by the product-induced topology;
its kernel consists exactly of representatives belonging to both
$N_0$ and $N_{\rm C}$, which is $K$. The definition of the quotient
topology therefore proves continuity and injectivity, without asserting
that this injection is a topological embedding into the product.

Writing $\pi_0,\pi_{\rm C}$ for its two coordinate projections, their
images and kernels are exactly
\begin{align}
 \operatorname{im}\pi_0&=\{[f]_{N_0}:f\in C\},&
 \ker\pi_0&=(N_0\cap\mathcal H)/K,\notag\\
 \operatorname{im}\pi_{\rm C}&=\{[f]_{N_{\rm C}}:f\in C\},&
 \ker\pi_{\rm C}&=(\A\cap N_{\rm C})/K.
 \label{ccb:kernels}
\end{align}
For example $\pi_0([f]_K)=0$ means precisely $f\in N_0$, while
$f\in C$ already supplies $f\in\mathcal H$; this proves the first
kernel identity in both directions. The other identity follows by
the same direct test in the second coordinate. Formula
\eqref{ccb:kernels} gives an explicit representative rule for the
correspondence even when either coordinate has several possible lifts.

The first image in \eqref{ccb:kernels} is dense in $Q_0$.
To prove this, take $\phi\in\E$, and choose a fixed smooth cutoff
$\chi$ equal to one on $[-1,1]$ and zero outside $[-2,2]$.
Set $\phi_R(v)=\chi(v/R)\phi(v)$ for $R\geq1$.
Leibniz's formula bounds every seminorm of $\phi-\phi_R$ by a
finite sum of weighted integrals of derivatives of $\phi$ over
$|v|\geq R$, multiplied by bounded constants and factors $R^{-k}$.
Those integrals tend to zero by their absolute integrability. Thus
$\phi_R\to\phi$ in $\E$. Moreover
\[
 |T\phi_R(s)|\leq4\|\phi_R\|_{L^1}e^{2R|\Im s|},
\]
so $T\phi_R$ is of order at most one and belongs to $C$.
Consequently $C$ is dense in $\A$; taking images under the continuous
surjective quotient map proves the assertion. All generators
\eqref{ccb:qf} also lie in $C$ by their proved $\E$ representatives
and their order at most one.

Finally this same correspondence transports exactly through every
complex Newman time. For $t\in\C$, define
\begin{equation}
 j_t:\mathscr R\longrightarrow Q_t\times Q_{\rm C},\qquad
 j_t([f]_K)=([U_tf]_{N_t},[f]_{N_{\rm C}}).
 \label{ccb:heat-correspondence}
\end{equation}
It is continuous and injective: $U_t$ is a continuous automorphism
and $U_tf\in N_t$ is equivalent to $f\in N_0$. Its two kernels as
coordinate projections are the same as in \eqref{ccb:kernels}, and its
first image is dense in $Q_t$. These assertions follow by applying
the proved isomorphism $Q_0\to Q_t$ to the first coordinate.
The space $C$ is stable under $M_s$: it is stable in $\A$ by
integration by parts and in $\mathcal H$ because the latter is the
specified topological ring. Each closed relation subspace is stable
under $M_s$, by continuity and invariance of $\C[s]\Xi$.
Thus $[f]_K\mapsto[sf]_K$ is defined, continuous, and obeys
\begin{equation}
 j_t([sf]_K)=
 \bigl([B_tU_tf]_{N_t},[sf]_{N_{\rm C}}\bigr),
 \qquad B_t=M_s-(t/2)\partial_s.
 \label{ccb:operator-correspondence}
\end{equation}
This is the exact operator correspondence: the first coordinate has
the heat-transported multiplication operator, while the second has
the original Connes multiplication operator. It follows directly from
$U_tM_s=B_tU_t$ and retains the full derivative term. No inclusion of
all of $\A$ into $\mathcal H$, no equality of their closures, and no
identification of their spectra is needed for the proved maps
\eqref{ccb:quotient-correspondence}--\eqref{ccb:operator-correspondence}.
